\documentclass[11pt]{article}

\usepackage[margin=1in]{geometry}
\usepackage{graphicx}
\usepackage{float}
\usepackage{amsmath}
\usepackage{amssymb}
\usepackage{booktabs}
\usepackage{xcolor}
\usepackage{hyperref}
\usepackage{enumitem}
\usepackage{caption}

\hypersetup{colorlinks=true, linkcolor=blue, urlcolor=blue, citecolor=blue}

\title{\textbf{Progress Update}}
\date{\today}

\begin{document}
\maketitle

\section{Input}

Two new recordings, exported directly through your own \texttt{ros2\_unbag}
pipeline rather than our \texttt{.mcap} fallback -- both include native
registered depth alongside RGB, gripper, and wrench, confirming the
compressedDepth export fix is working end to end.

\textbf{Recording 5.} Gripper, wrench, RGB, depth all present. 65.5\,s,
975 RGB frames. On inspection this recording contains three distinct
interaction phases in one continuous capture: a latch/hinge contact, a
Rubik's cube pick-press-drop cycle, then a second latch contact.

\textbf{Recording 6.} Same sensor profile. Short 9.7\,s clip, 156 frames.
Latch/hinge contact only, no grasp cycle.

\section{Experiments}

\begin{enumerate}[noitemsep]
  \item Full pipeline (event detection, seed placement, SAM\,2
        propagation, sidecar) on both recordings.
  \item Latch/hinge bolt hardware propagated as a \texttt{contact\_receiver}
        for the first time -- new object-role type, not seen in earlier
        recordings.
  \item Label-free object proposal (DINOv2 attention $\times$ real depth)
        scored quantitatively against the propagated ground-truth object,
        extended from one recording to every recording with real depth
        (16 events across 6 recordings total).
  \item Point-prompt vs.\ box-prompt SAM\,2 seeding compared by
        propagating both across a full 975-frame recording.
  \item Both candidate camera-bracket transforms from the CAD drawings
        (sent earlier) tested directly by running the live projection
        pipeline and measuring where the projected wrench ray lands.
  \item A direct hand-eye calibration procedure (standard method: printed
        calibration board, several arm poses, closed-form solve) written
        and validated against synthetic ground truth.
\end{enumerate}

\section{Results}

\textbf{Multi-phase event detection.} A single-event detector that reports
only the largest force peak in the whole recording silently attributes
the wrong phase's contact to Recording 5's cube task (the latch phase's
peak is larger). \textbf{Mitigation:} restricting the search to the
window between the detected grasp and release events recovers the correct
cube-related event; already the default elsewhere in the pipeline, now
applied here too.

\textbf{New object role.} Latch/hinge bolt hardware propagated cleanly as
a contact-receiver across Recording 6's full 156 frames.

\begin{figure}[H]
\centering
\includegraphics[width=0.6\textwidth]{figures/t001labexport/identify/overlays/f0079_1783420306137171476.png}
\caption{Latch bolt hardware (magenta), propagated as a contact-receiver
for the first time.}
\end{figure}

\textbf{Label-free discovery, quantified.} Mean intersection-over-union
between the label-free proposal and the true object is 0.161 across all
16 events (range 0.044--0.352) -- consistently low regardless of scene or
event type. \textbf{Mitigation:} none needed; extends the earlier
finding with a real number instead of a single-recording example.

\textbf{Point vs.\ box prompts.} Comparable while the object is genuinely
in frame. Once the object leaves frame, a box-prompted track degrades
gracefully toward empty (mean 3{,}243\,px); a point-prompted track does
not recognise absence and drifts onto unrelated background (mean
247{,}594\,px, 46\% of frame; peak 495{,}142\,px, 95\% of frame, on a bare
panel with no object visible). \textbf{Mitigation:} box prompts used for
any multi-coloured or irregular object going forward.

\textbf{CAD-based calibration.} Both candidate lens positions from the
bracket drawings (63\,mm apart, matching the ZED Mini's stereo baseline)
were tested directly: both project the wrench ray entirely outside the
visible frame, on every recording tested, regardless of which is chosen.
This indicates the CAD-derived bracket orientation is unreliable, not
only the ambiguous translation. \textbf{Mitigation:} CAD-reading is ruled
out as a calibration method here; a direct hand-eye calibration script
(printed board + several arm poses + closed-form solve) has been written
and validated against synthetic ground truth (recovers a known transform
exactly). Ready to run once there is a short window of rig access.

\section{Asks}

No blocking ask this update -- the earlier request to have someone extract
the bracket transform from CAD is no longer necessary given the result
above, so that effort can be dropped. When convenient, a short recording
of the arm pausing at 10--15 distinct poses in front of a printed
calibration board would let us close the calibration loop directly; happy
to send the exact board spec whenever useful.

\end{document}
